% SEGMENT OLP-0348-S01
% Part: lambda-calculus
% Chapter: introduction
% Section: lambda-definability

\documentclass[../../../include/open-logic-section]{subfiles}

\begin{document}

\olfileid{lam}{int}{rep}
\olsection{\usetoken{S}{lambda definable} எண்கணிதச் சார்புகள்}

லாம்டா கலனம் எவ்வாறு ஒரு கணக்கீட்டு மாதிரியாகப் பயன்பட முடியும்? முதலில்,
இக்கூற்றை எவ்வாறு பொருள்கொள்வது என்பதே தெளிவாக இல்லை. இயல் எண்கள் மீதான
கணிப்புத்தன்மையைப் பற்றிப் பேச, அவ்வெண்களுக்குப் பொருத்தமான ஒரு பிரதிநிதித்துவம்
தேவை. வியக்கத்தக்க அளவுக்கு நன்றாகச் செயல்படும் ஒன்று இதோ.

\begin{defn}
ஒவ்வோர் இயல் எண்~$n$ க்கும், மொத்தம் $n$ $x$-கள் உள்ள
$\lambd[x][\lambd[y][(x(x(x(\dots x(y)))))]]$ என்ற லாம்டா சொல்லை
\emph{சர்ச் எண்ணுரு}~$\num{n}$ என வரையறுக்க.
\end{defn}

சொற்கள் $\num{n}$ ``மீள்செயல்படுத்திகள்'': உள்ளீடு $f$ மீது, $y$ ஐ
$f^n(y)$ க்கு அனுப்பும் சார்பை $\num{n}$ திருப்பித்தருகிறது. ஒவ்வோர் எண்ணுருவும்
இயல்புவடிவில் உள்ளதைக் கவனிக்க. இயல் எண்கள் மீதான ஒரு சார்பை ஒரு லாம்டா சொல்
``கணிக்கிறது'' என்பதன் பொருளை இப்போது கூறலாம்.

\begin{defn}
% SOURCE CORRECTION TA-LAM-004: the displayed arguments use k, so the function is k-ary rather than n-ary.
$f(x_0, \dots, x_{k-1})$ ஆனது $\Nat$ இலிருந்து $\Nat$ குச் செல்லும் ஒரு
$k$-இட பகுதிச் சார்பாக இருக்கட்டும். ஒரு $\lambd$-சொல்~$F$ ஆனது~$f$ ஐ
\emph{!!{lambda define}s} என்று கூறுவது, அப்போதும் அப்போது மட்டுமே, இயல்
எண்களின் ஒவ்வொரு தொடர் $n_0$, \dots,~$n_{k-1}$ க்கும் பின்வருவன அமையும்.
$f(n_0, \dots, n_{k-1})$ வரையறுக்கப்பட்டால்,
\[
F\, \num{n_0}\, \num{n_1} \dots \num{n_{k-1}} \red \num{f(n_0, n_1, \dots,
  n_{k-1})}
\]
என்பது அமையும்.
% SOURCE CORRECTION TA-LAM-005: the undefined case applies F to the same numeral arguments; the frozen comma is not an argument separator.
அம்மதிப்பு வரையறுக்கப்படாவிட்டால், $F\, \num{n_0}\,
\num{n_1} \dots \num{n_{k-1}}$ க்கு இயல்புவடிவம் இல்லை.
\end{defn}

\begin{thm}
\ollabel{thm:lambda-def}
ஒரு சார்பு~$f$ பகுதியாகக் கணிக்கத்தக்கதாக இருப்பதும், அது ஒரு லாம்டா சொல்லால்
!!{lambda defined} ஆக இருப்பதும் ஒன்றுக்கொன்று நிகரானவை.
\end{thm}

\begin{explain}
இத்தேற்றம் ஓரளவு வியப்பூட்டுகிறது. கணக்கீட்டு மாதிரியாக லாம்டா கலனம் மிகவும்
எளிய கலனம்; லாம்டா அருவமாக்கமும் செயல்படுத்தலும் மட்டுமே அதன் செயல்கள்!{}
எனினும், இவ்வளவு குறைவான வளங்களிலிருந்தே எந்தக் கணக்கீட்டுச் செயல்முறையையும்
நடைமுறைப்படுத்த முடியும்.
\end{explain}

\end{document}
